\chapter{Subiecte examen}

\section{Examen AM2, 2017--2018}

\textbf{NR I}

1. Determinați funcția olomorfă $ f : \CC \to \CC, f = u + iv$, știind că:
\[
  \mathrm{Re}f = u(x, y) = x^2 - y^2 - 2y, \quad u : \RR^2 \to \RR
\]
și că $ f(1) = 1 $.

\vspace{1cm}

2. Să se calculeze integralele:

(a) Folosind teorema lui Cauchy:
\[
  \int_{|z| = 3} e^z + \frac{z + 8}{z-1} dz;
\]

(b) $ \ds\int_0^{2\pi} \frac{1}{(2 + \cos t)^2} dt $.

\vspace{1cm}

3. Fie funcția:
\[
  f : \CC - \{0, 2\} \to \CC, \quad f(z) = \frac{e^{2z}}{z^2(z - 2)}.
\]

Să se determine:
\begin{enumerate}[(a)]
\item $ \mathrm{Rez}(f, 0), \mathrm{Rez}(f, 2) $;
\item $ \ds\int_{|z - 2| = 3} f(z) dz $.
\end{enumerate}

\vspace{1cm}

4. Să se determine soluția ecuației diferențiale, folosind transformata Laplace:
\[
  x^{\prime\prime} - 4x^\prime + 3x = 3t, \quad x(0) = 2, x'(0) = 3.
\]

\vspace{1cm}

5. Să se rezolve ecuația integrală:
\[
  \varphi(t) = e^t + \int_0^t e^u \varphi(t - u) du.
\]

\vspace{2cm}

\textbf{NR 2.}

1. Să se determine funcția olomorfă $ f : \CC \to \CC, f = u + iv$, dacă:
\[
  \mathrm{Imf} = v(x, y) = x^3 - 3xy^2, \quad v : \RR \to \RR
\]
și știind că $ f(1) = 1 $.

\vspace{1cm}

2. Să se calculeze integralele:

(a) Folosind teorema lui Cauchy:
\[
  \int_{|z| = 3} z^2 + \frac{z + 7}{z - 2} dz.
\]

(b) $ \ds\int_{-\pi}^\pi \frac{\cos 3t}{5 - 4\cos t} dt $.

\vspace{1cm}

3. Fie funcția:
\[
  f : \CC - \{1, 3\} \to \CC, f(z) = \frac{e^z}{(z - 1)(z - 3)^2}.
\]
Să se determine:
\begin{enumerate}[(a)]
\item $ \mathrm{Rez}(f, 1), \mathrm{Rez}(f, 3) $;
\item $ \ds\int_{|z - 3| = 4} f(z) dz $.
\end{enumerate}

\vspace{1cm}

4. Să se rezolve ecuația diferențială, folosind transformata Laplace:
\[
  x^{\prime\prime} - 6x^\prime + 5x = 4e^t, \quad x(0) = 2, x'(0) = -1.
\]

\vspace{1cm}

5. Să se rezolve ecuația integrală:
\[
  \varphi(t) = t + 4 \int_0^t \frac{t - u}\varphi(u) du.
\]

%%%%%%%%%%%%%%%%%%%%%%%%%%%%%%%%%%%%%%%%%%%%%%%%%%%%%%%%%%%%%%%%%%%%%% 

\section{Restanță AM2, 2017--2018}

1. Rezolvați ecuațiile diferențiale:
\begin{enumerate}[(a)]
\item $ y^{\prime\prime} - 3y' + 2y = x + 1 $;
\item $ x^2 y^{\prime\prime} - 5xy' + 9y = x, x > 0 $.
\end{enumerate}

\vspace{1cm}

2. Să se determine liniile de cîmp pentru:
\[
  \vec{V}(x, y, z) = yz\vec{i} + xz\vec{j} + (x + y)\vec{k}.
\]

\vspace{1cm}

3. Să se aducă la forma canonică și să se rezolve ecuația:
\[
  \begin{cases}
    4 \frac{\partial^2 u}{\partial x^2} &= \frac{\partial^2 u}{\partial t^2} \\
    u(x, 0) &= x + 2 \\
    \frac{\partial u}{\partial t}(x, 0) &= 3x^2 + 1.
  \end{cases}
\]

\vspace{1cm}

4. Să se determine funcția olomorfă $ f : \CC \to \CC, f = u + iv $ știind
că $ v = x^4 + 6x^2y^2 + y^4, f(1) = 1 $.

\vspace{1cm}

5. Fie funcția complexă:
\[
  f : \CC - \{-1, 1 \} \to \CC, \quad %
  f(z) = \frac{e^z}{(z + 1)^2(z - 1)}.
\]
\begin{enumerate}[(a)]
\item Să se calculeze $ \mathrm{Rez}(f, 1), \mathrm{Rez}(f, -1) $;
\item Să se calculeze $ \ds\int_{|z + 1| = 3} f(z) dz $.
\end{enumerate}

\vspace{1cm}

%%%%%%%%%%%%%%%%%%%%%%%%%%%%%%%%%%%%%%%%%%%%%%%%%%%%%%%%%%%%%%%%%%%%%% 

6. Să se rezolve ecuația $ x^{\prime\prime} - 2x' - 5x = 16t \cdot e^t $,
știind că $ x(0) = 3, x'(0) = 4 $.

\vspace{1cm}

%%%%%%%%%%%%%%%%%%%%%%%%%%%%%%%%%%%%%%%%%%%%%%%%%%%%%%%%%%%%%%%%%%%%%% 

7. Calculați integrala:
\[
  \int_{|z + 1| = 4} \Big( z^2 + \frac{z + 7}{z - 2} \Big) dz.
\]

\vspace{1cm}

%%%%%%%%%%%%%%%%%%%%%%%%%%%%%%%%%%%%%%%%%%%%%%%%%%%%%%%%%%%%%%%%%%%%%% 

8. Calculați $ \ds\int_0^{2\pi} \frac{1}{2 + \cos \theta} d\theta $.


%%% Local Variables:
%%% mode: latex
%%% TeX-master: "../am2-18-19"
%%% End:
